% Lösungen Bedingte Wahrscheinlichkeit und Bayes (zusammenbau v0.4, Heft)
\einheitenkopf{Bedingte Wahrscheinlichkeit und Bayes}

\erg{11}{a) die Fußball spielenden Mitglieder \quad b) $\frac{0{,}16}{0{,}4} = 0{,}4$ \quad c) $\frac{36}{90} = 0{,}4$}
\erg{12}{a) $\frac{0{,}14}{0{,}4} = 0{,}35$ (geteilt durch den Anteil der Frauen, nicht durch $0{,}35$) \quad b) Kinderstation: $\frac{0{,}04}{0{,}1} = 0{,}4$; übrige: $\frac{0{,}15}{0{,}9} \approx 0{,}167$; der Anteil ist auf der Kinderstation größer (nicht $4\,\%$ mit $15\,\%$ vergleichen)}
\erg{13}{a) mit dem Baum \quad b) $\frac{0{,}1 \cdot 0{,}8}{0{,}1 \cdot 0{,}8 + 0{,}9 \cdot 0{,}1} = \frac{0{,}08}{0{,}17} \approx 0{,}471$ \quad c) die Wahrscheinlichkeit, dass ein defektes Gerät aus dem ersten Werk stammt: im Zähler der Pfad „erstes Werk und defekt“, im Nenner die Summe beider Pfade zu „defekt“}
\erg{14}{a) $\frac{0{,}6 \cdot 0{,}09}{0{,}4 \cdot 0{,}06 + 0{,}6 \cdot 0{,}09} = \frac{0{,}054}{0{,}078} \approx 0{,}692$ (im Nenner beide Pfade zu „unzufrieden“) \quad b) die Wahrscheinlichkeit, dass ein positiv getestetes Tier infiziert ist: Zähler der Pfad „infiziert und positiv“, Nenner alle Pfade zu „positiv“ (nicht: positiv, wenn infiziert) \quad c) $\frac{0{,}09}{0{,}06 + 0{,}09} = 0{,}6$, also ja (nicht mit dem Anteil unter den Aufträgen ohne A vergleichen) \quad d) Ja-Pfade: $0{,}9 \cdot 0{,}15 = 0{,}135$ (verlassen) und $0{,}1 \cdot 0{,}85 = 0{,}085$ (bleiben); $\frac{0{,}135}{0{,}22} \approx 0{,}614$ (nicht alle $15\,\%$ in den Zähler: wer bleiben will und die zweite Frage hat, sagt auch Ja)}
\erg{15}{a) nur im Nenner \quad b) $P = \frac{0{,}6a}{0{,}6a + 0{,}1 \cdot (1 - a)}$; Probe mit $a = 0{,}4$: $P = 0{,}8$ \quad c) Im Term steht b nur im Nenner, beim Pfad der reklamierenden Stammkunden; der Zähler, der Pfad der reklamierenden Neukunden, bleibt gleich. Sinkt b, sinkt der Nenner, also wächst der Bruch}
\erg{16}{a) Graph B: Kommt jedes andere Paket pünktlich an, stammt jedes verspätete von D, der Wert ist eins; kommen alle anderen zu spät, bleibt ein Wert über null, weil auch Pakete von D zu spät kommen – A endet bei null, C beginnt unter eins \quad b) $x \approx 0{,}15$; x ist der Rückgang des Anteils mit hohem Verbrauch unter den Fernwärme-Haushalten; $f(x) = 0{,}6$ heißt: 60\,\% der Haushalte mit hohem Verbrauch nutzen Fernwärme (nicht: 60\,\% der Fernwärme-Haushalte verbrauchen viel)}
